% Section 3
\section{Identification in \texorpdfstring{R$\times$C}{RxC} under linear contextual effects}\label{sec:rxc}

This section concerns the pooled model M0 and the superpopulation target $B^*$.
Throughout, assume $\E[N]<\infty$ and $\mu_r>0$ for every $r$.

\subsection{The unit polytope and what the data identify}

\begin{lemma}[The unit polytope]\label{lem:polytope}\status{proved}
For all $x\in\Delta^{R-1}$ and $t\in\Delta^{C-1}$: (a) $S(x,t)\neq\emptyset$;
(b) $S(x,t)$ is compact and convex; (c) $(x,t)\mapsto S(x,t)$ has a closed graph,
is upper hemicontinuous and Borel measurable; (d) $h_S(u;x,t)$ is sublinear and
$\sqrt R$-Lipschitz in $u$, concave in $t$, and Borel in $(u,x,t)$; (e)
$S(x,t)=\{b:\langle u,b\rangle\le h_S(u;x,t)\ \forall u\}$.
\end{lemma}
\begin{reading}
Each unit's set of possible rate tables is never empty (``every group behaves
like the unit as a whole'', $b_{rc}=t_c$, always fits) and it is a well-behaved
convex polytope. These are the regularity facts every later proof uses.
\end{reading}

\begin{lemma}[Identification decomposition]\label{lem:decomp}\status{proved}
Let $\mathbb Q\in\mathcal M(H;A)$ have $W$-marginal $P$, and assume A5. Then
(a) $\E[T_c\mid X=x]=x^\top Q_cx$ on $A$, and $Q_c$ is unique;
(b) $\sum_cH_c=\one\one^\top$, so $\sum_cQ_c=\one\one^\top$ and $\sum_cK_c=0$;
(c) every $K'\in\Kspace$ gives the same regression function, and
$\dim\Kspace=(C-1)R(R-1)/2$.
\end{lemma}
\begin{reading}
The data reveal the symmetric part $Q$ of the contextual-effect matrices and
nothing about the skew part $K$. For each pair of groups and each outcome, the
data identify the sum ``effect of the $s$-share on group $r$ plus effect of the
$r$-share on group $s$'', but not how it splits.
\end{reading}
\begin{sketch}
Conditional expectations of the accounting identity give $\E[T_c\mid X]=X^\top
H_cX$, and a skew matrix vanishes in a quadratic form, so only $Q_c$ appears. A5
makes the quadratic form determine $Q_c$. Row sums of $\beta$ equal one; under A5
this forces $\sum_cH_c=\one\one^\top$.
\end{sketch}

\begin{remark}[A5 in practice]\label{rem:A5}
For $R=2$, A5 holds iff $\supp_AP_X$ contains at least three distinct points. In
general A5 says that the design of the quadratic regression has full rank. For
$R=3$ that needs six support points not on a common conic, and any five points
lie on a conic. Fitting each column by ordinary least squares (OLS) on the
$R(R+1)/2$ quadratic monomials gives $\sum_c\hat Q_c=\one\one^\top$ exactly,
because $\sum_cT_{ic}=1=X_i^\top\one\one^\top X_i$; no constrained regression is
needed.
\end{remark}

\begin{lemma}[The estimand]\label{lem:estimand}\status{proved}
Under (S1)--(S2) with $P(X\in A)=1$,
\[
B^*_{rc}=\frac1{\mu_r}\Big[\sum_sQ_{c,rs}\mu_{rs}+\sum_{s\ne r}K_{c,rs}\mu_{rs}\Big],
\]
which is affine in $K$, with coefficient $\mu_{rs}/\mu_r\ge0$ on the entry $K_{c,rs}$.
\end{lemma}
\begin{reading}
The target is a known number plus a weighted sum of the unknown entries of $K$.
The weight $\mu_{rs}/\mu_r$ measures how much group $r$ lives alongside group
$s$: the more mixed the units, the more the unknown split matters. In pure
units it does not matter at all.
\end{reading}

\subsection{The sharp identified set}

Let $P_{T\mid V}$ be a regular conditional distribution of $T$ given $V$, and
\[
\psi(u;v):=\int h_S(u;x,t)\,P_{T\mid V=v}(dt),
\]
the conditional expectation of the support function at $V=v=(x,n)$.

\begin{theorem}[Characterisation of the sharp set]\label{thm:sharp}\status{proved}
Assume A1, A5 and $H=Q(P)+K$ with $K\in\Kspace$. Then $K\in\ThI(P;A)$ if and
only if (S3) holds and there is one $P_V$-null set $\mathcal N$ such that for
every $v=(x,n)\notin\mathcal N$ with $x\in A$,
\begin{equation}\label{eq:T1}
\langle u,M_H(x)\rangle\le\psi(u;v)\qquad\text{for all }u\in\R^{R\times C}.
\end{equation}
$\ThI(P;A)$ is convex and closed. It is compact when $\affdim(A)=R-1$, which A5
implies.
\end{theorem}
\begin{reading}
A value of $K$ is consistent with the data exactly when, at every composition and
size, the implied mean rate table could be the average of rate tables each of
which fits some unit's margins. Formally, $M_H(x)$ lies in the conditional
selection expectation of the random polytope $S(X,T)$ given $V=v$. This is the
strongest possible use of the data under the model. It needs the whole
conditional distribution of outcomes given composition and size, which in real
data has to be smoothed, so it is not what the package computes. Choosing $u$
and $-u$ along the equality constraints shows that \eqref{eq:T1} also requires
$\E[T\mid V]=\E[T\mid X]$, an observable implication of A3.
\end{reading}
\begin{sketch}
Necessity: each unit's $\beta$ lies in its polytope, so its conditional mean
satisfies the support-function inequality direction by direction. One null set
for all directions comes from a countable dense set of directions and the
Lipschitz continuity of $h_S$ in $u$. Sufficiency: the set of measurable
selections of $S(X,T)$ is weakly compact in $L^1$ (Dunford--Pettis and Mazur),
and so is its image under $\E[\cdot\mid V]$. If the target mean were outside that
image, the Hahn--Banach theorem would separate it by a bounded $V$-measurable
direction, and \eqref{eq:T1} evaluated at that random direction (a measurable
argmax attains the support function) would give a contradiction. No auxiliary
randomisation is needed, because $S(x,t)$ is convex. The random-set
formulation follows \citet{bmm2011}.
\end{sketch}

\begin{corollary}[Sharp bounds for a linear target]\label{cor:sharp}\status{proved}
If $P(X\in A)=1$ and $\ThI\neq\emptyset$, the sharp identified set for
$\varphi=\sum_{rc}a_{rc}B^*_{rc}$ is the interval
$[\min_{K\in\ThI}\varphi(K),\ \max_{K\in\ThI}\varphi(K)]$, with $\varphi(K)$ given
by Lemma~\ref{lem:estimand}.
\end{corollary}
\begin{reading}
The unknown enters the target linearly and the consistent set is convex and
compact, so the answer is always an interval.
\end{reading}

\paragraph{Continuous versus integer rates.} Sharpness is relative to real-valued
latent rates. With integer cell counts the exact feasible set is a subset of
$S(X,T)$. Every vertex of $S(x,t)$ with integer margins is an integer table,
because transportation polytopes are totally unimodular, so the Duncan--Davis
bounds are attained by integer tables. Whether randomising over integer vertices
leaves the sharp set unchanged has only a sketch of an argument, and nothing in
the paper uses it.

\subsection{The outer set: reduction, nesting and sharpness for the regression}

\begin{lemma}[Vertex reduction, boundedness, validity]\label{lem:outer}\status{proved}
(a) If $A$ is a polytope and $\ell$ is linear, $\ell\ge0$ on $A$ iff $\ell(v)\ge0$
at every vertex $v$ of $A$; so the constraints of $\ThO$ may be imposed at
$\vertx(A)$ only. (b) If $\affdim(A)=R-1$, $\ThO(P;A)$ is bounded, hence a
polytope; otherwise it is a closed convex polyhedron, possibly unbounded.
(c) Under (S1)--(S2), if $A\subseteq\conv(\supp_AP_X)$, then (S3) holds
automatically; in particular for $A=\conv\supp P_X$.
\end{lemma}
\begin{reading}
Only the most extreme units matter for the constraints. If the region $A$ is the
range of the data, probability validity of the straight lines is automatic. If
$A$ extrapolates beyond the data, or reaches into a part of the hull where the
model is imposed on no support point, validity there is an extra assumption.
For example, with $A$ a single interior point that carries no mass, (S2)
restricts nothing and $M_H$ need not be valid there; and a single-point $A$ can
make $\ThO$ unbounded. A5 rules out the second case, because it forces $A$ to be
full-dimensional.
\end{reading}

\begin{proposition}[Nesting]\label{prop:nesting}\status{proved}
$\ThI(P;A)\subseteq\ThX(P;A)\subseteq\ThO(P;A)$.
\end{proposition}
\begin{reading}
Using more of the data can only narrow the set. The outer set uses only the
regression function, so it is the widest of the three. The second inclusion is
JKST's ``Jensen'' step; in R$\times$C it needs no Jensen inequality, because
conditional expectations of valid rate tables are valid rate tables.
\end{reading}

\begin{proposition}[$\ThO$ is sharp for the regression function]\label{prop:sharpreg}\status{proved}
Let $\mathcal P(P)$ be the set of laws $P'$ of $(N,X,T')$ with the same law of
$(N,X)$ as $P$ and $\E_{P'}[T'\mid V]=\E_P[T\mid X]$. If $P(X\in A)=1$, then
\[
\ThO(P;A)=\bigcup_{P'\in\mathcal P(P)}\ThI(P';A).
\]
\end{proposition}
\begin{reading}
The outer bound is exactly what can be learned from the regression alone. It is
the sharp set for any data set in which units with the same composition have
identical outcomes. Calling it ``sharp'' without that qualifier would be wrong.
\end{reading}
\begin{sketch}
Given $K\in\ThO$, set every unit's rates equal to the implied mean, $\beta=M_H(X)$.
These are valid rates because $K\in\ThO$, they reproduce the regression, and $K$
is in the sharp set of the resulting law. The reverse inclusion follows from
Proposition~\ref{prop:nesting}, since $\ThO$ depends on $P$ only through $Q(P)$.
\end{sketch}

\begin{proposition}[When the outer set is sharp]\label{prop:IequalsO}\status{proved}
Assume $P(X\in A)=1$. (a) If $T=\E[T\mid X]$ almost surely, then $\ThI=\ThO$.
(b) More generally, suppose $\E[T\mid V]=\E[T\mid X]$ and, for $P_V$-almost every
$v=(x,n)$ and every $u$, $t\mapsto h_S(u;x,t)$ is affine on the convex hull of the
conditional support of $T$ given $V=v$. Then $\ThI=\ThO$. The condition in (b) is
sufficient, not necessary.
\end{proposition}
\begin{reading}
The regression-only bound loses nothing unless units with similar compositions
are so different that some sit at a Duncan--Davis corner and others do not. In
$2\times2$, the condition says that among units with the same composition the
outcome share never lies on both sides of $x$ or of $1-x$, where the unit bounds
have their kinks. In 400 random valid $2\times2$ designs the inclusion
$\ThX\subseteq\ThO$ was strict in 46.
\end{reading}

\subsection{Closed forms and the linear program}

\begin{proposition}[Equivalence with JKST]\label{prop:jkst}\status{proved}
Let $R=C=2$, $X$ the row-1 share, $A=\{(t,1-t):t\in[l,u]\}$ with $0<l<u<1$, and
write $\E[T_1\mid X]=w_0+c_1X+d_1X^2$ and $K=K_{1,12}$. Then (a)
$Q_{1,22}=w_0$, $Q_{1,12}=w_0+c_1/2$, $Q_{1,11}=w_0+c_1+d_1$ and $K=c_1/2-w_1$,
where $w_1$ is JKST's unidentified slope; (b) $\ThO$, written in terms of $w_1$, is
JKST's Proposition~2; (c) with $A=\Delta^1$ it is JKST's Proposition~3; (d) if
$P(l\le X\le u)=1$, $B^*(K)$ is the population version of JKST's equation~(16).
\end{proposition}
\begin{reading}
The R$\times$C machinery contains the published $2\times2$ method exactly, and
the package's original bounds are a special case. Numerically, the homogeneous
form reproduces the original implementation to $10^{-16}$. The resulting bounds
are over the outer set, as are JKST's: sharp for the regression function, not for
the full law.
\end{reading}

\begin{proposition}[The whole simplex]\label{prop:simplex}\status{proved}
With $A=\Delta^{R-1}$, $\ThO=\{K\in\Kspace:\ |K_{c,rs}|\le Q_{c,rs}\ \forall c,\ r<s\}$,
which also requires $Q_{c,rr}\ge0$ for all $c,r$ (a testable condition not involving $K$).
It is non-empty iff every $Q_{c,rs}\ge0$; then each $K_{c,rs}$ ranges over
$[-m_{c,rs},m_{c,rs}]$ with $m_{c,rs}=\min\{Q_{c,rs},1-Q_{c,rs}\}\le\tfrac12$, and
\[
B^*_{rc}\in\frac1{\mu_r}\Big[\sum_sQ_{c,rs}\mu_{rs}\mp\sum_{s\ne r}m_{c,rs}\mu_{rs}\Big].
\]
A negative $Q_{c,rs}$ refutes the model with $A=\Delta^{R-1}$.
\end{proposition}
\begin{reading}
If the straight lines are trusted all the way to pure units, two hypothetical
probabilities with a known average $Q$ can split it by at most $\pm\min(Q,1-Q)$.
This is narrower than the hull-based bound but assumes extrapolation.
\end{reading}

\begin{lemma}[Box intersected with a hyperplane]\label{lem:box}\status{proved}
Let $\ell\le h$ in $\R^C$ and $\mathcal P=\{\kappa:\ell\le\kappa\le h,\ \sum_c\kappa_c=0\}$.
(a) $\mathcal P\neq\emptyset$ iff $\sum_c\ell_c\le0\le\sum_ch_c$.
(b) The projection of $\mathcal P$ on coordinate $c$ is
$\big[\max\{\ell_c,-\sum_{c'\ne c}h_{c'}\},\ \min\{h_c,-\sum_{c'\ne c}\ell_{c'}\}\big]$.
(c) $\max_{\mathcal P}g^\top\kappa$ is attained by the greedy rule: start from
$\kappa=\ell$ and raise coordinates in decreasing order of $g_c$, each up to
$h_c$, until $\sum_c\kappa_c=0$.
\end{lemma}

\begin{proposition}[Two rows, any number of columns]\label{prop:tworows}\status{proved}
Let $R=2$ and $\vertx(A)=\{(t_k,1-t_k)\}_{k=1,2}$ with $0<t_1<t_2<1$. Define
\[
\lambda_c(t)=-Q_{c,12}-Q_{c,11}\tfrac{t}{1-t},\qquad
\rho_c(t)=Q_{c,12}+Q_{c,22}\tfrac{1-t}{t},\qquad
\ell_c=\max_k\lambda_c(t_k),\qquad h_c=\min_k\rho_c(t_k).
\]
Then $\ThO=\{\kappa:\ell\le\kappa\le h,\ \sum_c\kappa_c=0\}$, each $\kappa_c$ is
projected by Lemma~\ref{lem:box}(b), and if $P(X\in A)=1$ the bounds on
$B^*_{1c}=[Q_{c,11}\mu_{11}+(Q_{c,12}+\kappa_c)\mu_{12}]/\mu_1$ (increasing in
$\kappa_c$), on $B^*_{2c}=[(Q_{c,12}-\kappa_c)\mu_{12}+Q_{c,22}\mu_{22}]/\mu_2$
(decreasing in $\kappa_c$), and on any linear target (by the greedy rule) are in
closed form.
\end{proposition}
\begin{reading}
For two groups and any number of outcomes, each outcome's unknown is squeezed
between the tightest lower and upper limits implied by the two extreme units (the
least and the most group-1-heavy), and the unknowns must sum to zero across
outcomes. No LP solver is needed. Group 1's rate rises and group 2's rate falls
as the unknown rises, because the unknown is the share of the observed pattern
attributed to group 1.
\end{reading}

\begin{corollary}[Joint versus collapsed analysis]\label{cor:joint}\status{proved}
Let $R=2$. The collapsed analysis of column $c$ treats the table as two columns,
$c$ and ``not $c$''. (a) The joint interval for $\kappa_c$ lies inside the
collapsed one. (b) They are equal whenever, for each row $r$, the numbers
$Q_{c',rr}$ ($c'\neq c$) share one sign; in particular when every extrapolated
pure-unit rate $H_{c',rr}=Q_{c',rr}$ is non-negative. (c) Strict tightening
therefore requires some negative extrapolated pure-unit rate, so the linear model
is not valid on the whole simplex.
\end{corollary}
\begin{reading}
With two groups, fitting all outcome columns jointly helps a single column only if
the straight lines would give an impossible rate at some pure unit. With three or
more groups the rows are coupled through the shared entries of $K$, and joint
fitting can be strictly tighter even when every extrapolation is valid (strict
in 20 of the $3\times3$ cases checked). A characterisation for $R\ge3$ is open.
Joint fitting is needed anyway for targets that combine columns.
\end{reading}

\begin{proposition}[The general linear program]\label{prop:lp}\status{proved}
Assume $P(X\in A)=1$ and that $A$ is a polytope. In the free coordinates $k$,
with $K_c=\sum_{r<s}k_{c,rs}(e_re_s^\top-e_se_r^\top)$ for $c<C$ and
$K_C=-\sum_{c<C}K_c$,
\[
\ThO=\{k:\ Gk\le g\},
\]
with one row for each (vertex $v$, row $r$, column $c$): $-e_r^\top K_c(k)v\le(Q_cv)_r$.
A linear target is $\varphi(k)=\varphi_0+\omega^\top k$ with
\[
\varphi_0=\sum_{rc}\frac{a_{rc}}{\mu_r}\sum_sQ_{c,rs}\mu_{rs},\qquad
\omega_{c,rs}=\mu_{rs}\Big[\frac{a_{rc}-a_{rC}}{\mu_r}-\frac{a_{sc}-a_{sC}}{\mu_s}\Big]\quad(c<C,\ r<s),
\]
and its bounds are $\varphi_0+\min/\max\{\omega^\top k:Gk\le g\}$. With $A$ the
sample hull, imposing the rows at every observed $X_i$ is equivalent to imposing
them at the hull's vertices.
\end{proposition}
\begin{reading}
The computable bound is a small LP: one variable per unknown direction and one
constraint per extreme unit, group and outcome. It is an LP for the outer set,
not an LP form of the sharp set, which in general needs infinitely many
conditional constraints.
\end{reading}

\paragraph{Regularity of the value function.} Write
$V(Q,\mu)=\max\{\omega(\mu)^\top k:Gk\le g(Q)\}$. For fixed $\mu$, $V$ is concave
and piecewise linear in $Q$ where the LP is feasible and bounded. For fixed $Q$
it is convex and piecewise linear in $\omega$. Jointly it is in general neither:
$\max\{\mu k:-Q\le k\le Q\}=\mu Q$ is bilinear. $V$ is differentiable where the
optimal basis is unique and non-degenerate, which generalises JKST's
tie-breaking condition. Section~\ref{sec:inference} conditions on the design, so
$\omega$ and $G$ are fixed and only the first case is needed.

\begin{proposition}[Units outside the region]\label{prop:trim}\status{proved}
Let $A$ be a polytope with $P(X\in A)>0$, (S1) everywhere, (S2)--(S3) on $A$,
and A5 on $\supp_AP_X$. Let $\varphi^A_0$ and $\omega^A$ be the quantities of
Proposition~\ref{prop:lp} with $\mu_{rs}$ replaced by $\mu^A_{rs}=\E[NX_rX_s\one\{X\in A\}]$,
$\tilde a(x)_{rc}=a_{rc}x_r/\mu_r$, and
\[
U_{\rm out}=\E\big[N\,h_S(\tilde a(X);X,T)\one\{X\notin A\}\big],\qquad
L_{\rm out}=-\E\big[N\,h_S(-\tilde a(X);X,T)\one\{X\notin A\}\big].
\]
Then $\varphi\in\big[\varphi^A_0+\min_{\ThO}\omega^{A\top}k+L_{\rm out},\
\varphi^A_0+\max_{\ThO}\omega^{A\top}k+U_{\rm out}\big]$, and each endpoint is
attained. For a single cell the outside terms are the Duncan--Davis terms of the
units outside $A$.
\end{proposition}
\begin{reading}
Units outside the region where straight lines are trusted contribute their
model-free range; nothing links them to the model. This is what the package's
\code{range = "trim"} option computes for single cells.
\end{reading}

\subsection{Refutability and undetectable violations}

\begin{proposition}[Non-emptiness and refutation]\label{prop:nonempty}\status{proved}
If $P$ is generated by some $\mathbb Q\in\mathcal M(H_0;A)$ with $P(X\in A)=1$,
then the true $K_0\in\ThI\subseteq\ThO$ and
$B^*_{rc}\in[\varphi_L,\varphi_U]\cap\DD^*_{rc}$, where
\[
\DD^*_{rc}=\Big[\frac{\E[N\max\{0,X_r+T_c-1\}]}{\mu_r},\ \frac{\E[N\min\{X_r,T_c\}]}{\mu_r}\Big].
\]
An empty $\ThO$, or an empty intersection with $\DD^*$, refutes A1--A5 jointly.
\end{proposition}
\begin{reading}
If the model is right, the population bound is never empty. An empty bound is
evidence against the model; in samples it can also be noise, and
Section~\ref{sec:inference} turns it into a test with controlled level.
\end{reading}

\begin{proposition}[What the data can never rule out]\label{prop:undetectable}\status{proved}
Let $\nu:\Delta^{R-1}\to\R^{R\times C}$ be Borel with, at every $x\in\supp P_X$,
(i) $x^\top\nu(x)=0$, (ii) $\nu(x)\one=0$, and (iii) $M_H(x)+\nu(x)\in\Bset$. Then
the structures $\beta=M_H(X)+\nu(X)$ and $\beta'=M_H(X)$ generate the same law of
the data, but
$B^*_{rc}(\beta)-B^*_{rc}(\beta')=\E[NX_r\nu_{rc}(X)]/\mu_r$. Under A5 on $\supp_AP_X$, with $A\supseteq\supp P_X$, the linear
such $\nu$ are exactly $M_{K'}$ with $K'\in\Kspace$.
\end{proposition}
\begin{reading}
There are nonlinear contextual effects that cancel in the aggregate and leave the
data identical, yet change the transition rates. For two rows, $\nu_{1c}=b_cx_1x_2$
and $\nu_{2c}=-b_cx_1^2$ (with $\sum_cb_c=0$) give a quadratic contextual effect
for row 1 that is exactly cancelled in row 2; JKST's Example~6 is the case $C=2$.
The linearity assumption is therefore a genuine leap that no fit statistic can
confirm. Section~\ref{sec:sensitivity} bounds its consequences with a
sensitivity parameter.
\end{reading}

\subsection{Illustration}

In a noise-free $3\times3$ design with 2000 units, rows and columns ordered
(in-party, out-party, non-voter) and an arbitrary linear data-generating process,
the outer bounds on four transition rates are narrower than Duncan--Davis and
contain the truth:

\begin{center}\small
\begin{tabular}{@{}lcccc@{}}
\toprule
Target & Truth & $\ThO$ bound & Duncan--Davis & Width ratio \\
\midrule
Loyalty P$\to$P & 0.848 & [0.396, 0.932] & [0.131, 0.973] & 0.64 \\
Defection P$\to$O & 0.086 & [0.022, 0.361] & [0.000, 0.628] & 0.54 \\
Demobilisation P$\to$N & 0.066 & [0.022, 0.301] & [0.000, 0.532] & 0.52 \\
Mobilisation N$\to$P & 0.169 & [0.048, 0.590] & [0.000, 0.896] & 0.60 \\
\bottomrule
\end{tabular}
\end{center}

The bounds remain wide because six numbers are unidentified. Narrowing them is
the job of the district rungs and of the optional restrictions in the next two
sections. One unknown, for example, is $k_{P,PO}$, half the gap between the
loyalty of in-party voters in a unit otherwise all out-party and the rate at
which out-party voters switch to the in-party in a unit otherwise all in-party;
the data identify the average of these two rates, not the gap.

\paragraph{Relation to Cross and Manski.} In the framework of
\citet{crossmanski2002}, the data-alone region for a unit is the polytope
$S(x,t)$, and their instrument assumption with no further covariates is
Goodman's constancy, the special case of LCE with zero contextual effects,
$H_c=\alpha_c\one^\top$. Their outer region from iterated expectations
corresponds to $\ThO$, and their sharp region to $\ThI$. LCE is weaker than
constancy: rates may vary linearly with composition, so the result is a set
rather than a point. It is also a mean restriction on a superpopulation rather
than an equality of each unit's rates, which is why the bounds concern
aggregates, why each unit stays at its Duncan--Davis polytope, and why the gap
between $B$ and $\Bt$ must be handled in inference.
